\chapter{Chronic Low Back Pain}
\index{Chronic Low Back Pain}
\label{sec:Chronic Low Back Pain}

\section{Overview}
Chronic low back pain (cLBP) is pain in the lower back persisting for $\ge$3 months. It is the leading cause of years lived with disability worldwide, with a lifetime prevalence of 60--80\%. Most cases are non-specific (no identifiable structural pathology).
\section{Causes}
This condition happens because of a biopsychosocial interaction of biological factors (structural pathology, inflammation), psychological factors (catastrophizing, fear-avoidance beliefs, depression), and social factors (occupational demands, disability compensation).     
\section{Risk Factors}

\begin{itemize}[noitemsep]
  \item Genetic predisposition and family history
  \item Advancing age
  \item Lifestyle factors (diet, exercise, smoking, alcohol)
  \item Environmental exposures
  \item Underlying medical conditions
  \item Medication use
\end{itemize}
\section{Signs and Symptoms}

\subsection{Common symptoms}
\begin{itemize}[noitemsep]
  \item \item You may experience pain in the lower back that may radiate to the legs (sciatica), with or without neurological deficits. Pain or discomfort in the affected area    \item Pain or discomfort in the affected area
  \end{itemize}

\subsection{Emergency warning signs}
\begin{itemize}[noitemsep]
  \item Severe or worsening symptoms of chronic low back pain require immediate medical evaluation
\end{itemize}
\section{Progression and Stages}
The condition may progress through different stages of severity over time. Early stages often involve mild or subtle symptoms that may be overlooked. Moderate stages involve more noticeable symptoms affecting daily function. Advanced stages may involve significant impairment and complications.
\section{Complications}
\begin{itemize}[noitemsep]
  \item Disease progression leading to worsening symptoms
  \item Secondary infections or injuries
  \item Side effects of treatment
  \item Reduced quality of life
  \item Emotional and psychological impact
\end{itemize}
\section{Diagnosis}
\begin{itemize}[noitemsep]
  \item Doctors usually diagnose this based on your symptoms and a physical exam
  \item Red flags for serious pathology include age >50 or <20, recent trauma, unexplained weight loss, fever, night pain, history of cancer, and progressive neurological deficits. Initial Treatment options include education, activity modification, NSAIDs, and muscle relaxants
  \item First-line pharmacotherapy is NSAIDs, with TCAs, SNRIs, and gabapentinoids as adjuncts
  \item Non-pharmacological treatments include physical therapy, manual therapy, acupuncture, thinking and memory behavioral therapy, and graded activity programs
  \item Interventional procedures include epidural steroid injections, facet joint injections, and spinal cord stimulation.
  \item Comprehensive medical history and physical examination
  \item Laboratory tests to assess organ function
  \item Imaging studies to visualize affected structures
  \item Specialized testing depending on the affected system
  \item Consultation with relevant specialists
\end{itemize}
\section{Prevention}
\begin{itemize}[noitemsep]
  \item Maintain a healthy lifestyle with balanced diet and exercise
  \item Avoid known risk factors
  \item Attend regular medical check-ups for early detection
  \item Follow recommended screening guidelines
\end{itemize}
\section{Treatment Options}
Medications to manage symptoms and address underlying mechanisms Lifestyle modifications and self-care measures Physical therapy and rehabilitation Surgical intervention when indicated Regular monitoring and follow-up care Patient education and support services

\begin{itemize}[noitemsep]
  \item Medications to manage symptoms and address underlying mechanisms
  \item Lifestyle modifications and self-care measures
  \item Physical therapy and rehabilitation
  \item Surgical intervention when indicated
  \item Regular monitoring and follow-up care
\end{itemize}
\section{Living with It}
\begin{itemize}[noitemsep]
  \item Follow your treatment plan as prescribed
  \item Maintain regular follow-up with your healthcare provider
  \item Make necessary lifestyle adjustments
  \item Seek support from family, friends, or support groups
  \item Stay informed about your condition
\end{itemize}
\section{Outlook}
\begin{itemize}[noitemsep]
  \item The outlook varies from person to person
  \item Most cases improve with conservative management.
\end{itemize}
